
\clearpage
\chapter{עבודות קשורות: אבטחת \textenglish{AI} ושרשרת אספקה}

\section{איומים על מודלי \textenglish{LLM}}

מחקרים מוקדמים התמקדו בהתקפות \textenglish{adversarial examples} על מודלי \textenglish{NLP} . עם הופעת \textenglish{ChatGPT} וסוכנים, מוקד האיומים עבר: \cite{perez2022ignore} הדגים כי ניתן "לאפס" \textenglish{LLM} עם הוראות פשוטות טמונות בפרומפט, עוקף כל \textenglish{system prompt}.

\section{הזרקת פרומפט עקיפה}

\cite{greshake2023indirect} הציג תקיפה מתוחכמת יותר: הוראות נסתרות בתוכן שה-\textenglish{LLM} \textit{קורא} (אתרי אינטרנט, מסמכי \textenglish{PDF}, תוצאות \textenglish{API}). המחקר הראה פגיעות ב-\textenglish{7} מתוך \textenglish{7} מודלים שנבדקו. \textenglish{SkillSieve} מתמודד עם תקיפה זו דרך ניתוח \textenglish{sandbox} שחוסם גישה לרשת בזמן הבדיקה.

\section{הגנות קיימות}

\cite{liu2023prompt} סקר \textenglish{13} שיטות הגנה כנגד הזרקת פרומפט. שיטות מבוססות \textenglish{input sanitization} נכשלות מול קידודים יצירתיים; שיטות מבוססות \textenglish{fine-tuning} יעילות יותר אך יקרות. \textenglish{SkillSieve} נקט גישה כלאיים (\textenglish{hybrid}) שאינה תלויה באימון ספציפי להתקפה.

\section{אבטחת שרשרת אספקה תוכנה}

בעולם התוכנה המסורתי, \textenglish{SolarWinds} (2020) ו-\textenglish{XZ Utils} (2024) הדגימו את חומרת התקפות שרשרת האספקה. מנגנוני \textenglish{SLSA} (\textenglish{Supply chain Levels for Software Artifacts}) ו-\textenglish{Sigstore} נועדו לפתור בעיות אלה אך טרם הותאמו לפצ'קות (\textenglish{skill packages}) של מסגרות \textenglish{AI} \cite{weidinger2021ethical}.

\section{מצב האמנות}

\textenglish{SkillSieve} ממצב את עצמו כפתרון ייעודי לצינורות \textenglish{LLM}-סוכן, מעבר לפתרונות כלליים כמו \textenglish{Snyk} ו-\textenglish{OWASP Dependency-Check} שאינם מיועדים לניתוח סמנטי של תיאורי כישורים \cite{nakash2024skillsieve}.
